\begin{figure}[H]
\centering
\includegraphics[width=1.0\linewidth]{fidelity/G7-01-hokkaido_trajectory_map.pdf}
\caption[冬の 1 例（北海道）]{2020 年 1 月 15 日 00 UTC に VOISING から上げた場合の飛行経路。
左が全体、右が着地点の拡大。
この例では 4 通りの差が \SI{2.9}{\kilo\metre} しかなく、
和歌山の同じ日（図\ref{fig:g7-01-wakayama}、\SI{13.1}{\kilo\metre}）よりずっと小さい。
和歌山では基準（黒）と時刻のみ固定（青）、
地点のみ固定（橙）と両方固定（黄）がそれぞれ対になるのに対し、
ここでは 4 本が順に並ぶだけで、そうした対応は見えない。
\textbf{1 例は 1 例であって統計ではない}。
どちらが典型かは表\ref{tab:fidelity-error} で判断してほしい。}
\label{fig:g7-01-hokkaido}
\end{figure}

\begin{figure}[H]
\centering
\includegraphics[width=1.0\linewidth]{fidelity/G7-07-hokkaido_trajectory_map.pdf}
\caption[夏の 1 例（北海道）]{同じ図の 7 月版。
拡大パネルの縮尺が冬と夏で違う点に注意。}
\label{fig:g7-07-hokkaido}
\end{figure}